\documentclass[11pt]{article}

\usepackage[utf8]{inputenc}
\usepackage[T1]{fontenc}
\usepackage{amsmath,amssymb,amsfonts,bm,booktabs}
\usepackage{geometry}
\geometry{margin=1in}

\title{Supplementary Material:\\
MacroIRT and MacroMRT --- Bayesian Updates with State-Dependent Noise}
\author{}
\date{}

\begin{document}
\maketitle

\newcommand{\E}{\mathbb{E}}
\newcommand{\Var}{\operatorname{Var}}
\newcommand{\diag}{\operatorname{diag}}
\newcommand{\Nch}{N_{\mathrm{ch}}}
\newcommand{\gbarz}{\overline{\boldsymbol{\gamma}}_0}
\newcommand{\sbarz}{\overline{\boldsymbol{\sigma}^2}_0}
\newcommand{\Gbarmat}{\overline{\boldsymbol{\Gamma}}}
\newcommand{\Vbarmat}{\overline{\boldsymbol{V}}}
\newcommand{\Vbarv}{\overline{\boldsymbol{V}}_{\!v}}
\newcommand{\tildegg}{\widetilde{\gamma^{T}\Sigma\gamma}}
\newcommand{\tildevv}{\widetilde{v^{T}\Sigma v}}
\newcommand{\tildegvec}{\widetilde{\gamma^{T}\Sigma}}
\newcommand{\tildevvec}{\widetilde{v^{T}\Sigma}}
\newcommand{\Sprop}{\boldsymbol{\Sigma}^{\mathrm{prop}}}
\newcommand{\Spost}{\boldsymbol{\Sigma}^{\mathrm{post}}}
\newcommand{\mupost}{\boldsymbol{\mu}^{\mathrm{post}}}
\newcommand{\muprior}{\boldsymbol{\mu}^{\mathrm{prior}}}
\newcommand{\SmD}{\mathrm{SmD}}

\begin{abstract}
This supplement derives MacroIRT (interval-averaged Bayesian update with
state-dependent open-channel noise) and its M-family counterpart MacroMRT,
in the rank-1 quasi-Laplace formulation. The derivation follows the
``deepseek'' presentation: expansion at $p=\muprior$, Sherman--Morrison
inversion, boundary lift via the tilde operator. The full rank-2 exact
Laplace (which keeps the $\tfrac12\log V$ term and gives a $2\times2$
Woodbury inversion) is preserved in
\texttt{theory/macroir/archive/Macro\_IRT/macroirt\_macroirt\_supplement.tex}
for reference; for typical $\Nch\gtrsim 100$ it makes a negligible
difference.
\end{abstract}

% ============================================================
\section{Setup}
% ============================================================

We have $\Nch$ independent Markov ion channels with $K$ microscopic states.
At the start of an interval $[0,t]$, the macroscopic occupancy distribution is
$\mathbf{r} \sim \mathcal{N}\!\bigl(\boldsymbol{\mu}_0,\ \boldsymbol{\Sigma}_0/\Nch\bigr)$
(Moffatt 2007 scaling: $\boldsymbol{\Sigma}_0$ is the per-channel covariance,
$\Nch$-independent). The microscopic dynamics are governed by the rate
matrix $\mathbf{Q}$ and transition matrix
$\mathbf{P}(t) = \exp(\mathbf{Q}t)$.

Per boundary pair $(i_0,i_t)$, the model defines two boundary-conditioned
moments:
\begin{itemize}
  \item $\Gbarmat(i_0,i_t)$: boundary mean current
        $= \E[\,\overline{I}_{0\to t}\,|\,X(0)=i_0, X(t)=i_t\,]$.
        \quad[code: \texttt{gmean\_ij}]
  \item $\Vbarmat(i_0,i_t)$: boundary residual variance
        $= \Var[\,\overline{I}_{0\to t}\,|\,X(0)=i_0, X(t)=i_t\,]$.
        \quad[code: \texttt{gvar\_ij}]
\end{itemize}

Their per-$i_0$ marginals (the M-family inputs):
\begin{align}
  \gbarz(i_0)
  &= \sum_{i_t} P_{i_0\to i_t}(t)\,\Gbarmat(i_0,i_t)
   \quad\text{[code: \texttt{gmean\_i}]}\\
  \sbarz(i_0)
  &= \sum_{i_t} P_{i_0\to i_t}(t)\,\Vbarmat(i_0,i_t)
   \quad\text{[code: \texttt{gvar\_i\_residual} (computed on the fly per the in-place fix)]}
\end{align}

Centred prior covariance: $\SmD := \boldsymbol{\Sigma}_0 - \diag(\boldsymbol{\mu}_0)$.
Instrument noise: $e := \epsilon^2/t + \nu^2$.

\paragraph{Important: residual vs. total $\sbarz$.}
$\sbarz$ above is the \emph{residual} per-$i_0$ variance: the part of
$\Var(\overline{I}\,|\,X_0)$ that is NOT explained by where the trajectory
ends ($X_t$). The variance-of-conditional-mean across end states is
\emph{separately} captured by the boundary cross-correlation term inside
$\tildegg$ (Section~\ref{sec:predictive}). Using the residual $\sbarz$
here, with $\tildegg$ separately, avoids the double-count diagnosed in
\texttt{theory/macroir/notes/gvar\_i\_overcount\_audit.md}.

% ============================================================
\section{Observation model and predictive moments}
\label{sec:predictive}
% ============================================================

The interval-averaged observation $\overline{y}^{\mathrm{obs}}$, given the
macroscopic occupancy $\mathbf{r}$, is heteroscedastic Gaussian:
\begin{equation}
  \overline{y}^{\mathrm{obs}}\,\big|\,\mathbf{r}
  \,\sim\,
  \mathcal{N}\bigl(\Nch\,\mathbf{r}\!\cdot\!\boldsymbol{\gamma},\
                   e + \Nch\,\mathbf{r}\!\cdot\!\boldsymbol{\sigma}^2\bigr),
\end{equation}
where for now we treat $\boldsymbol{\gamma}$ and $\boldsymbol{\sigma}^2$ as
abstract per-state vectors; the boundary lift in
Section~\ref{sec:boundary-lift} replaces them with their boundary-conditioned
analogues.

At $\mathbf{r} = \boldsymbol{\mu}_0$ we obtain the \emph{predictive mean and
variance} (the actual quantities the Kalman update consumes):
\begin{align}
  \overline{y}^{\mathrm{pred}}_{0\to t}
  &= \Nch\,\boldsymbol{\mu}_0\cdot\gbarz,\\
  V
  &= e \,+\, \Nch\,\tildegg \,+\, \Nch\,\boldsymbol{\mu}_0\cdot\sbarz,
  \label{eq:V}
\end{align}
where $\tildegg$ is the boundary cross-correlation term (defined in
Section~\ref{sec:boundary-lift}); for the M-family (no boundary lift) it
collapses to $\gbarz^{T}\,\boldsymbol{\Sigma}_0\,\gbarz$.

Innovation:
\begin{equation}
  \delta := \overline{y}^{\mathrm{obs}} - \overline{y}^{\mathrm{pred}}_{0\to t}.
\end{equation}

% ============================================================
\section{Energy, gradient, Hessian (rank-1 quasi-Laplace)}
% ============================================================

We work with the negative log-posterior \emph{quasi-energy} (dropping the
$\tfrac12\log V$ term that gives the rank-2 correction):
\begin{equation}
  E(\mathbf{r})
  \,=\,
  \frac{(y - \Nch\,\mathbf{r}\!\cdot\!\boldsymbol{\gamma})^2}
       {2(e + \Nch\,\mathbf{r}\!\cdot\!\boldsymbol{\sigma}^2)}
  \,+\,
  \tfrac{\Nch}{2}\,(\mathbf{r}-\muprior)\,\boldsymbol{\Sigma}_p^{-1}\,(\mathbf{r}-\muprior)^{T}.
\end{equation}

Differentiating at $\mathbf{r} = \muprior$ (so the prior gradient vanishes):
\begin{equation}
  \mathbf{g}
  \,=\,
  -\,\frac{\Nch\,\delta}{V}\,\boldsymbol{\gamma}
  \,-\,
  \frac{\Nch\,\delta^2}{2V^2}\,\boldsymbol{\sigma}^2.
  \label{eq:grad_explicit}
\end{equation}
Define the \emph{tilted direction}
\begin{equation}
  \boxed{\
  \mathbf{v} \,:=\, \boldsymbol{\gamma} + \frac{\delta}{V}\,\boldsymbol{\sigma}^2
  \quad\text{[code: \texttt{v}]}
  \ }
  \label{eq:vdef}
\end{equation}
which factors the gradient into a single rank-1 direction:
\begin{equation}
  \mathbf{g}
  \,=\,
  -\,\frac{\Nch\,\delta}{2V}\bigl(\boldsymbol{\gamma} + \mathbf{v}\bigr).
  \label{eq:grad_factored}
\end{equation}

The Hessian is rank-1 in the $\mathbf{v}$ direction (the second-derivative
chain rule absorbs all four $(\boldsymbol{\gamma},\boldsymbol{\sigma}^2)$
outer products into $\mathbf{v}\mathbf{v}^{T}$):
\begin{equation}
  \mathbf{H}
  \,=\,
  \Nch\,\boldsymbol{\Sigma}_p^{-1}
  \,+\,
  \frac{\Nch^2}{V}\,\mathbf{v}\mathbf{v}^{T}.
  \label{eq:Hessian}
\end{equation}

\paragraph{Why the gradient is rank-2 in $(\boldsymbol{\gamma},\boldsymbol{\sigma}^2)$
but the Hessian is rank-1 in $\mathbf{v}$.}
The two terms in \eqref{eq:grad_explicit} use different powers of $\delta$
($\delta^1$ vs $\delta^2$). The Hessian's mixed derivatives recombine these
into a single $\mathbf{v}\mathbf{v}^{T}$ outer product because
$\partial(\mathbf{v}\mathbf{v}^{T})/\partial\mathbf{r}$ at $\mathbf{r}=\muprior$
generates exactly the cross-derivative terms. This is the structural reason
the rank-1 form is so clean: \emph{the heteroscedasticity is fully absorbed
into a single direction $\mathbf{v}$}.

% ============================================================
\section{Sherman--Morrison inversion}
% ============================================================

Define the scalar
\begin{equation}
  \boxed{\
  s \,:=\, \mathbf{v}^{T}\,\boldsymbol{\Sigma}_p\,\mathbf{v}
  \quad\text{[code: \texttt{s}]}
  \ }
\end{equation}
(in the M-family this is just the quadratic form; in the I-family it
becomes the tilde scalar $\tildevv$ of Section~\ref{sec:boundary-lift}).

Sherman--Morrison applied to \eqref{eq:Hessian}:
\begin{equation}
  \mathbf{H}^{-1}
  \,=\,
  \frac{1}{\Nch}\,\boldsymbol{\Sigma}_p
  \,-\,
  \frac{1}{V + \Nch\,s}\,\boldsymbol{\Sigma}_p\,\mathbf{v}\mathbf{v}^{T}\,\boldsymbol{\Sigma}_p.
\end{equation}

The (Moffatt-scaled) posterior covariance is $\Spost := \Nch\,\mathbf{H}^{-1}$:
\begin{equation}
  \boxed{\
  \Spost
  \,=\,
  \boldsymbol{\Sigma}_p
  \,-\,
  \frac{\Nch}{V + \Nch\,s}\,\boldsymbol{\Sigma}_p\,\mathbf{v}\mathbf{v}^{T}\,\boldsymbol{\Sigma}_p.
  \ }
  \label{eq:Sigma_post_MRT}
\end{equation}

Posterior mean via Newton step from $\mathbf{r}_0 = \muprior$:
\begin{equation}
  \mupost
  \,=\,
  \muprior \,-\, \mathbf{H}^{-1}\mathbf{g}
  \,=\,
  \muprior \,-\, (\Spost/\Nch)\,\mathbf{g}.
\end{equation}
With $\mathbf{g}$ from \eqref{eq:grad_factored}:
\begin{equation}
  \boxed{\
  \mupost
  \,=\,
  \muprior \,+\, \frac{\delta}{2V}\,\Spost\,(\boldsymbol{\gamma} + \mathbf{v}).
  \ }
  \label{eq:mu_post_MRT}
\end{equation}

% ============================================================
\section{Boundary-state lift via the tilde operator}
\label{sec:boundary-lift}
% ============================================================

For the I-family (interval-averaged measurement), $\boldsymbol{\gamma}$ and
$\boldsymbol{\sigma}^2$ in the formulas above are not per-state constants but
boundary-conditioned moments $\Gbarmat$ and $\Vbarmat$. The tilde operator
expresses each of the four bilinear forms appearing in
\eqref{eq:Sigma_post_MRT}--\eqref{eq:mu_post_MRT} in terms of
$\boldsymbol{\Sigma}_p$ ($K\!\times\! K$), $\mathbf{P}$ ($K\!\times\! K$), and
$\Gbarmat,\Vbarmat$ ($K\!\times\! K$). \textbf{No $K^2\!\times\! K^2$ object is
ever materialised.}

\paragraph{Tilde scalar.}
For boundary moment matrices $\overline{\mathbf{U}},\overline{\mathbf{W}}$ with
per-$i_0$ marginals $\overline{\mathbf{u}}_0,\overline{\mathbf{w}}_0$:
\begin{equation}
  \widetilde{u^{T}\Sigma w}
  \,:=\,
  \overline{\mathbf{u}}_0^{T}\,\SmD\,\overline{\mathbf{w}}_0
  \,+\,
  \boldsymbol{\mu}_0^{T}\!\bigl[\,\mathbf{P}\circ\overline{\mathbf{U}}\circ\overline{\mathbf{W}}\,\bigr]\mathbf{1}.
\end{equation}

\paragraph{Tilde vector (endpoint frame).}
\begin{equation}
  \widetilde{u^{T}\Sigma}
  \,:=\,
  \overline{\mathbf{u}}_0^{T}\,\SmD\,\mathbf{P}
  \,+\,
  \boldsymbol{\mu}_0\,\bigl[\mathbf{P}\circ\overline{\mathbf{U}}\bigr].
\end{equation}

\paragraph{The two new tilde objects needed for IRT.}
The $\mathbf{v}$-direction lift requires the boundary-lifted analog of
$\mathbf{v}$ at the boundary level:
\begin{equation}
  \boxed{\
  \Vbarv \,:=\, \Gbarmat \,+\, \frac{\delta}{V}\,\Vbarmat
  \quad\text{[code: \texttt{gtotal\_v\_ij} (new), mirrors \texttt{gtotal\_ij} family]}
  \ }
\end{equation}
(elementwise on the $K\!\times\! K$ boundary matrices). Then:
\begin{align}
  \tildevvec
  &= \mathbf{v}^{T}\,\SmD\,\mathbf{P}
   + \boldsymbol{\mu}_0\,(\mathbf{P}\circ\Vbarv),
   &\text{[code: \texttt{vS} (new), mirrors \texttt{gS}]}
   \label{eq:tilde_v_vec}\\
  \tildevv
  &= \mathbf{v}^{T}\,\SmD\,\mathbf{v}
   + \boldsymbol{\mu}_0^{T}\!\bigl[\mathbf{P}\circ\Vbarv\circ\Vbarv\bigr]\mathbf{1}.
   &\text{[code: \texttt{vSv} (new), mirrors \texttt{gSg}]}
   \label{eq:tilde_v_v}
\end{align}
The $\boldsymbol{\gamma}$-direction lift (already in current MacroIR code):
\begin{align}
  \tildegvec
  &= \gbarz^{T}\,\SmD\,\mathbf{P}
   + \boldsymbol{\mu}_0\,(\mathbf{P}\circ\Gbarmat),
   &\text{[code: \texttt{gS}]}\\
  \tildegg
  &= \gbarz^{T}\,\SmD\,\gbarz
   + \boldsymbol{\mu}_0^{T}\!\bigl[\mathbf{P}\circ\Gbarmat\circ\Gbarmat\bigr]\mathbf{1}.
   &\text{[code: \texttt{gSg}]}
\end{align}

% ============================================================
\section{MacroIRT update equations}
% ============================================================

The IRT mean update is \emph{not} a direct ``substitute $(\boldsymbol{\gamma}+\mathbf{v}) \to (\tildegvec + \tildevvec)$ and keep $\Spost$ on the right'' of the M-family boxed equation. That naive substitution under-shoots by a factor $V/V_{\mathrm{pred}}$ at $\mathbf{v}=\boldsymbol{\gamma}$ and fails to reduce to MacroIR. The correct lift propagates the start-frame Newton step through $\mathbf{P}(t)$ and expands the Sherman--Morrison down-date \emph{before} applying the tilde substitution.

\paragraph{Derivation.}
The M-family Newton step at $p_0=\muprior$ in start frame is
\(
  \Delta p_{\text{start}}^{\text{col}}
  = (\delta/(2V))\,\Sigma_p^{\text{post}}(\boldsymbol{\gamma}+\mathbf{v})
\)
with the rank-1 SM-derived
\(
  \Sigma_p^{\text{post}} = (1/\Nch)\Sigma_p - (1/(V+\Nch s))\Sigma_p\mathbf{v}\mathbf{v}^{T}\Sigma_p
\),
$s=\mathbf{v}^{T}\Sigma_p\mathbf{v}$, $V=\epsilon^2+\Nch\,\boldsymbol{\mu}\!\cdot\!\boldsymbol{\sigma}^2$
(the energy-derivation value of $V$, equal to $V_{\mathrm{obs}}$ in code, \emph{not} the predictive $V_{\mathrm{pred}}=V+\Nch\tildegg$).
Projecting to endpoint frame as a row, $\Delta p_{\text{end}} = (\Delta p_{\text{start}}^{\text{col}})^{T}\mathbf{P}$, and expanding $\Sigma_p^{\text{post}}$:
\begin{align*}
  \Delta p_{\text{end}}
  &= (\delta/(2V))\,(\boldsymbol{\gamma}+\mathbf{v})^{T}\Sigma_p^{\text{post}}\,\mathbf{P}\\
  &= (\delta/(2V))\Bigl[(\boldsymbol{\gamma}+\mathbf{v})^{T}\Sigma_p\mathbf{P}
       \;-\; \frac{\Nch}{V+\Nch s}\,(\boldsymbol{\gamma}+\mathbf{v})^{T}\Sigma_p\mathbf{v}\;\mathbf{v}^{T}\Sigma_p\mathbf{P}\Bigr].
\end{align*}
The tilde lift then replaces each underlined-by-shape object by its boundary-conditioned analog --- as a single contracted unit, not as ``$\boldsymbol{\gamma}^{T}$ alone, with $\Sigma$ outside'':
\begin{itemize}
  \item $\boldsymbol{\gamma}^{T}\Sigma_p\mathbf{P} \;\to\; \tildegvec$
  \item $\mathbf{v}^{T}\Sigma_p\mathbf{P} \;\to\; \tildevvec$
  \item $\mathbf{v}^{T}\Sigma_p\mathbf{v} \;\to\; \tildevv$
  \item $\boldsymbol{\gamma}^{T}\Sigma_p\mathbf{v} \;\to\; b' := \widetilde{\gamma^{T}\Sigma v} = \tildegg + \alpha\,(\delta/V)\,b_{\text{tilde}}$,
        where $b_{\text{tilde}} := \widetilde{\gamma^{T}\Sigma\sigma^2}$ is the cross tilde scalar between the $\boldsymbol{\gamma}$ and $\boldsymbol{\sigma}^2$ directions.
\end{itemize}

\paragraph{Propagation.}
\begin{align}
  \muprior &= \boldsymbol{\mu}_0\,\mathbf{P}(t),
   &\text{[code: \texttt{p\_P\_mean() * t\_P()}]}\\
  \Sprop
  &= \mathbf{P}(t)^{T}\,\SmD\,\mathbf{P}(t) + \diag(\muprior),
   &\text{[code: \texttt{sigma\_pre}]}
\end{align}

\paragraph{Posterior covariance (rank-1 down-date, endpoint frame).}
\begin{equation}
  \boxed{\
  \Spost
  \,=\,
  \Sprop
  \,-\,
  \frac{\Nch}{V + \Nch\,\tildevv}\,(\tildevvec)^{T}\,\tildevvec.
  \ }
  \label{eq:Sigma_post_IRT}
\end{equation}
This rank-1 down-date \emph{is} correct as stated --- only the mean update needed the SM expansion.

\paragraph{Posterior mean (single Newton step from $\muprior$, corrected form).}
With $\mathrm{sm} := \Nch / (V + \Nch\,\tildevv)$ the SM factor used in $\Spost$,
\begin{equation}
  \boxed{\
  \mupost
  \,=\,
  \muprior
  \,+\,
  \frac{\delta}{2V}\,\Bigl[\,(\tildegvec + \tildevvec)\;-\;\mathrm{sm}\,(b' + \tildevv)\,\tildevvec\,\Bigr].
  \ }
  \label{eq:mu_post_IRT}
\end{equation}
This is the rigorous IRT mean update. The bracket can be read as ``ideal gradient direction $(\tildegvec+\tildevvec)$ minus the SM down-projection along the $\tildevvec$ direction''.

\paragraph{Reduction to MacroIR at $\boldsymbol{\sigma}^2 = 0$.}
With $\mathbf{v}=\boldsymbol{\gamma}$, $\tildevvec=\tildegvec$, $\tildevv=\tildegg$, $b'=\tildegg$, the bracket becomes
\(
  2\tildegvec - \mathrm{sm}\,(2\tildegg)\,\tildegvec = 2\tildegvec(1 - \Nch\tildegg/V_{\mathrm{pred}}) = 2\tildegvec\,V/V_{\mathrm{pred}}
\)
(using $V_{\mathrm{pred}}=V+\Nch\tildegg$), so $\mupost = \muprior + (\delta/V_{\mathrm{pred}})\,\tildegvec$ --- the canonical MacroIR endpoint Kalman update from \texttt{theory/macroir/docs/Macro\_IR/MacroIR\_tilde.md} \S4.2. \emph{Structural}, not coincidental. This was the test that the previous ``$(\tildegvec+\tildevvec)\Spost$'' form failed.

% ============================================================
\section{MacroMRT update equations}
% ============================================================

The M-family is the no-boundary-lift restriction. Replace each tilde
scalar/vector with its per-$i_0$ ensemble counterpart (drop the
$\boldsymbol{\mu}_0\,[\mathbf{P}\circ\dots]\mathbf{1}$ second term in
each tilde):
\begin{align}
  \tildegg &\to \gbarz^{T}\,\SmD\,\gbarz,\\
  \tildevv &\to \mathbf{v}^{T}\,\SmD\,\mathbf{v},\\
  \tildegvec &\to \gbarz^{T}\,\SmD\,\mathbf{P},\\
  \tildevvec &\to \mathbf{v}^{T}\,\SmD\,\mathbf{P}.
\end{align}
With these substitutions the same boxed update formulas apply. Importantly,
in the M-family the $\sbarz$ entering $\mathbf{v}=\gbarz + (\delta/V)\sbarz$
is the \emph{total} variance $\gbarz^{T}\!\diag\,\gbarz - \dots$ (or
equivalently $\mathtt{gsqr\_i} - \mathtt{gmean\_i}^2$), not the residual:
the M-family has no boundary cross-correlation in $\tildegg$, so the
total-variance form is what's needed (no double-count concern).

% ============================================================
\section{Reductions (the diamond)}
% ============================================================

\begin{center}
\begin{tabular}{l c c}
\toprule
& Boundary lift? & $\boldsymbol{\sigma}^2$ Taylor (i.e. $\mathbf{v}$ direction)? \\
\midrule
MacroMR  & no  & no \\
MacroMRT & no  & yes \\
MacroIR  & yes & no \\
MacroIRT & yes & yes \\
\bottomrule
\end{tabular}
\end{center}

\paragraph{IRT $\to$ IR (set $\sbarz = \mathbf{0}$).}
$\mathbf{v}=\gbarz$, $\Vbarv=\Gbarmat$, $\tildevvec=\tildegvec$,
$\tildevv=\tildegg$. The mean update collapses to
$\mupost = \muprior + (\delta/V)\,\tildegvec\,\Spost / 2 \cdot 2 = \muprior + (\delta/V)\,\tildegvec\,\Spost / \dots$
Actually more carefully: with $\boldsymbol{\gamma} + \mathbf{v} = 2\gbarz$ at
$\sbarz=0$, the factor $\delta/(2V) \cdot 2 = \delta/V$, recovering the
standard MacroIR update.

\paragraph{IRT $\to$ MRT (drop boundary lift).}
Per-$i_0$ substitutions of Section~7.

\paragraph{MRT $\to$ MR (set $\sbarz = \mathbf{0}$).}
Recovers Moffatt 2007.

% ============================================================
\section{Implementation map (math $\leftrightarrow$ code)}
% ============================================================

Naming policy: \textbf{keep all existing code names; introduce new names
only where no existing analog exists}. The four genuinely new objects
(rank-1 quasi-Laplace) are \texttt{v}, \texttt{gtotal\_v\_ij}, \texttt{vS},
\texttt{vSv}, named to mirror the existing \texttt{gtotal\_ij}/\texttt{gS}/\texttt{gSg}
patterns.

\begin{center}
\begin{tabular}{l l l}
\toprule
Math & Code name & Status \\
\midrule
$\boldsymbol{\mu}_0$ & \texttt{p\_P\_mean} & existing \\
$\boldsymbol{\Sigma}_0$ & \texttt{p\_P\_Cov} & existing \\
$\boldsymbol{\Sigma}_0 - \diag\boldsymbol{\mu}_0$ & \texttt{SmD} & existing \\
$\mathbf{P}(t)$ & \texttt{t\_P} & existing \\
$\Gbarmat$ (boundary mean current) & \texttt{gmean\_ij} & existing \\
$\Vbarmat$ (boundary residual var) & \texttt{gvar\_ij} & existing in \texttt{Qdt} \\
$\gbarz$ (per-$i_0$ mean) & \texttt{gmean\_i} & existing (also \texttt{t\_gmean\_i} in scope) \\
$\sbarz$ (per-$i_0$ residual var) & \texttt{gvar\_i\_residual} & added at \texttt{qmodel.h:4063} \\
$P_{ij}\,\Gbarmat$ & \texttt{gtotal\_ij} & existing \\
$P_{ij}\,\Vbarmat$ & \texttt{gtotal\_var\_ij} & existing in \texttt{Qdt} \\
$\delta$ (innovation) & \texttt{dy} & existing \\
$V$ (predictive variance) & \texttt{r\_y\_var} & existing \\
$\tildegvec$ (\(\boldsymbol{\gamma}\)-direction tilde vector) & \texttt{gS} & existing \\
$\tildegg$ (\(\boldsymbol{\gamma}\)-direction tilde scalar) & \texttt{gSg} & existing \\
$\muprior$ & \texttt{mu\_prior} (= \texttt{p\_P\_mean * t\_P}) & existing \\
$\Sprop$ & \texttt{sigma\_pre} & existing \\
\midrule
\multicolumn{3}{l}{\textit{New objects (rank-1 quasi-Laplace IRT/MRT):}}\\
$\mathbf{v} = \gbarz + (\delta/V)\sbarz$ & \texttt{v} & new \\
$\Vbarv = \Gbarmat + (\delta/V)\Vbarmat$ & \texttt{gtotal\_v\_ij} & new (mirrors \texttt{gtotal\_ij} family) \\
$\tildevvec$ & \texttt{vS} & new (mirrors \texttt{gS}) \\
$\tildevv$ & \texttt{vSv} & new (mirrors \texttt{gSg}) \\
$\mupost$ & \texttt{r\_P\_mean} & existing (replaces standard Kalman) \\
$\Spost$ & \texttt{r\_P\_cov} & existing (replaces standard Kalman) \\
\bottomrule
\end{tabular}
\end{center}

\paragraph{Note on alternative naming.}
The deepseek doc proposed math-faithful names (\texttt{gamma0\_bar},
\texttt{sigma2\_gamma0}, \texttt{Gamma\_bar}, \texttt{V\_bar}, \texttt{tilde\_v\_Sigma\_vec},
etc.). We do \emph{not} adopt these because they would parallel rather than
reuse the existing kernel naming. The existing names
(\texttt{gmean\_i}, \texttt{gmean\_ij}, \texttt{gvar\_ij}, \texttt{gS}, \texttt{gSg})
already follow a consistent pattern; the new objects (\texttt{v},
\texttt{gtotal\_v\_ij}, \texttt{vS}, \texttt{vSv}) extend that pattern naturally.

\paragraph{Already-dead names available for rank-2 reuse.}
The submission's orphaned Taylor block at \texttt{qmodel.h:3362} uses
\texttt{sS} (\(\sigma\)-direction tilde vector \(\widetilde{\sigma^{T}\Sigma}\)),
\texttt{sSs} (\(\widetilde{\sigma^{T}\Sigma\sigma}\)), \texttt{sSg}
(\(\widetilde{\sigma^{T}\Sigma\gamma}\)) for the \(\sigma^2\)-direction tilde
quantities. Rank-1 IRT/MRT does not need these; they are reserved for any
future rank-2 implementation.

\paragraph{Computational cost over MacroIR.}
For IRT, the new objects to compute per interval are: \texttt{v} ($O(K)$),
\texttt{gtotal\_v\_ij} ($O(K^2)$), \texttt{vS} ($O(K^2)$, mirrors existing
\texttt{gS} computation), \texttt{vSv} ($O(K^2)$, mirrors \texttt{gSg}),
the Sherman-Morrison factor ($O(1)$). Total overhead is roughly twice the
existing tilde computations, no new asymptotic complexity. For MRT, even
simpler: drop the boundary cross-correlation second term in each tilde
(\texttt{vSv} reduces to \texttt{TranspMult(v,sigma\_pre)*v}, etc.).

% ============================================================
\section{Pseudo-code}
% ============================================================

Throughout, \texttt{V\_obs = r\_y\_var - N*gSg} is the measurement-noise piece at $r=\muprior$ (the supplement's $V$ in the energy/Hessian derivation), \emph{not} the predictive \texttt{r\_y\_var = V\_obs + N*gSg}. Using \texttt{r\_y\_var} for $\beta = \delta/V$ in $\mathbf{v}$, and using \texttt{r\_y\_var} in the SM denominator of the mean step, were two of the original bugs.

\paragraph{IRT measurement update (av=2 + vc=true).}
\begin{verbatim}
// Inputs in scope from MacroIR setup:
//   p_P_mean, p_P_Cov, SmD, t_P, t_gmean_i (=gmean_i), gmean_ij, gvar_ij,
//   gtotal_ij, gtotal_var_ij, gSg, gS, mu_prior, sigma_pre, dy, r_y_var,
//   N (=N_ch), gvar_i_residual.

// 0. V_obs (the V in the energy/Hessian derivation; NOT r_y_var)
auto V_obs = r_y_var - N * gSg;

// 1. Effective direction (per-i_0 K-dim), β = δ/V_obs
auto v = t_gmean_i() + (dy / V_obs) * gvar_i_residual;

// 2. Boundary-lifted v_ij and gtotal_v_ij (used inside vSv and vS)
auto gtotal_v_ij = gtotal_ij + (dy / V_obs) * gtotal_var_ij;
auto v_ij        = gmean_ij  + (dy / V_obs) * gvar_ij;

// 3. Tilde scalars (mirrors gSg using the X̃ᵀΣỸ closed form
//    = TranspMult(X̄,SmD)·Ȳ + p·(gtotal_X_ij ∘ Y_ij)·u)
Matrix<double> u(p_P_mean().size(), 1, 1.0);
auto vSv     = getvalue(TranspMult(v, SmD) * v)
             + getvalue(p_P_mean() * (elemMult(gtotal_v_ij, v_ij) * u));
auto b_tilde = getvalue(TranspMult(t_gmean_i(), SmD) * gvar_i_residual)
             + getvalue(p_P_mean() * (elemMult(gtotal_ij, gvar_ij) * u));
auto b_prime = gSg + (dy / V_obs) * b_tilde;        // γ̃ᵀΣv

// 4. Tilde vector for v in endpoint frame
auto vS = TranspMult(v, SmD) * t_P() + p_P_mean() * gtotal_v_ij;

// 5. Sherman-Morrison rank-1 down-date (V_obs in denominator!)
double sm_factor = N / (V_obs + N * vSv);
auto Sigma_post  = sigma_pre - sm_factor * XTX(vS);

// 6. Mean update (single Newton step from mu_prior; row-vector convention;
//    the SM expansion is explicit, NOT (gS+vS)·Sigma_post)
auto mean_dir = (dy / (2 * V_obs)) *
                ((gS + vS) - sm_factor * (b_prime + vSv) * vS);
auto mu_post = mu_prior + mean_dir;

// 7. log L unchanged (uses predictive r_y_var, not V_obs)
double logL = -0.5 * (std::log(2 * M_PI * r_y_var) + dy * dy / r_y_var);
\end{verbatim}

\paragraph{MRT measurement update (av=1 + vc=true).}
Same skeleton as IRT but each tilde object collapses to its per-$i_0$
ensemble counterpart (drop the second $\boldsymbol{\mu}\,[\mathbf{P}\!\circ\!\dots]\mathbf{1}$ term).
The M-family uses the \emph{total} per-$i_0$ variance
\texttt{gvar\_i\_total = gsqr\_i - gmean\_i\(^2\)} (no double-count concern,
since the M-family has no boundary cross-cov in $\tildegg$):
\begin{verbatim}
auto V_obs = r_y_var - N * gSg;

// 1. Effective direction (M-family uses gvar_i_total, not residual)
auto v = t_gmean_i() + (dy / V_obs) * gvar_i_total;

// 3'. Scalar (M-family vSv collapses to start-frame quadratic)
double vSv     = getvalue(TranspMult(v, p_P_Cov) * v);
double b_tilde = getvalue(TranspMult(t_gmean_i(), p_P_Cov) * gvar_i_total);
double b_prime = gSg + (dy / V_obs) * b_tilde;       // γᵀΣ_p v

// 4'. Cross-cov vector vS = vᵀ·Σ_p·P (endpoint frame)
auto vS = TranspMult(v, p_P_Cov) * t_P;

// 5'. Sherman-Morrison rank-1 down-date
double sm_factor = N / (V_obs + N * vSv);
auto Sigma_post  = sigma_pre - sm_factor * XTX(vS);

// 6'. Mean update (same SM-expanded form)
auto mean_dir = (dy / (2 * V_obs)) *
                ((gS + vS) - sm_factor * (b_prime + vSv) * vS);
auto mu_post = mu_prior + mean_dir;
\end{verbatim}

% ============================================================
\section{Trust-region simplex shrink (preserved from MacroIR)}
% ============================================================

The existing trust-region machinery in
\texttt{safely\_calculate\_Algo\_State\_recursive} (lines 4108-4144 of
\texttt{legacy/qmodel.h}) ensures $\mupost$ stays inside the simplex floor
$[p_{\min},p_{\max}]$ and $\Spost$ stays PSD. The same machinery applies to
the IRT/MRT updates: treat the bracketed direction inside
\eqref{eq:mu_post_IRT} as the effective \texttt{chi}\,$\cdot$\,\texttt{gS} of the
standard rank-1 update, and the rank-1 outer product
$(\tildevvec)^{T}\tildevvec$ as the standard down-date direction. The
$\alpha_\mu$ and $\alpha_\sigma$ scalar minimisations apply unchanged.

% ============================================================
\section{Note on the rank-2 exact Laplace}
% ============================================================

The rank-1 quasi-Laplace formulation above drops the $\tfrac12\log V(\mathbf{r})$
term from the energy. The rank-2 \emph{exact} Laplace keeps it, which:
\begin{itemize}
  \item adds a second score component
        $\mathrm{score}_V = \delta^2/V - 1$ to the gradient,
  \item makes the Hessian rank-2 (with an additional
        $\boldsymbol{\sigma}^2\boldsymbol{\sigma}^{2T}$ outer product, negative-weighted),
  \item requires a $2\!\times\! 2$ Woodbury inversion instead of the scalar
        Sherman-Morrison,
  \item shifts $\mupost$ even when $\delta=0$ (toward states with smaller
        predicted variance).
\end{itemize}
For typical $\Nch\gtrsim 100$ this correction is $O(1/\Nch)$ relative to
the leading term; the rank-1 formulation is the operational form. The full
rank-2 derivation, written in the same canonical naming, is preserved in
\texttt{theory/macroir/archive/Macro\_IRT/macroirt\_macroirt\_supplement.tex}.
Adopt rank-2 only if small-$\Nch$ stress tests show the rank-1 form
breaking down.

\end{document}
